% from ega0-15-fr.tex
\oldpage[0\textsubscript{IV-1}]{20}
Soient $A$ un anneau, $M$ un $A$-module, $(f_i)_{1\leq i\leq n}$ une suite
$M$-régulière d'éléments de $A$. Alors, quels que soient les entiers
$\nu_i\geq1$, la suite $(f_i^{\nu_i})$ est $M$-régulière et
$M/(\sum_{i\leq n}f_i^{\nu_i}M)$ a une suite de composition dont les
quotients sont isomorphes à $M/(\sum_{i\leq n}f_iM)$.
\end{corollary}

\begin{proof}
Raisonnons par récurrence sur $n$~; pour $n=1$ les assertions sont immédiates
puisque $f^jM/f^{j+1}M$ est isomorphe à $M/fM$, l'homothétie de rapport $f$
dans $M$ étant injective par hypothèse. Posons alors
\[
  N=M/\left(\sum_{i\leq n-1}f_i^{\nu_i}M\right)~;
\]
par l'hypothèse de récurrence, $N$ admet une filtration finie $(N_k)$ pour
laquelle les $\gr_k(N)$ sont isomorphes à
$M/(\sum_{i\leq n-1}f_iM)$~; par hypothèse $f_n$ est donc
$\gr_\bullet(N)$-régulier~; il résulte par suite de (15.1.19) que $f_n$ est
$N$-régulier. Appliquant le cas $n=1$ démontré au début, on en conclut que
$f_n^{\nu_n}$ est $N$-régulier et que $N/f_n^{\nu_n}N$ admet une suite de
composition dont les quotients sont isomorphes à $N/f_nN$~; mais pour la
filtration de $N/f_nN$ quotient de la filtration $(N_k)$,
$\gr_k(N/f_nN)$ est isomorphe à $\gr_k(N)/f_n\gr_k(N)$ par (15.1.19) pour
tout $k$, donc à $M/(\sum_{i\leq n}f_iM)$, ce qui prouve le corollaire.
\end{proof}

\begin{proposition}[15.1.21]
\label{0.15.1.21-fr}
Soient $A$, $B$ deux anneaux locaux noethériens, $\varphi:A\to B$ un
homomorphisme local, $M$ un $B$-module de type fini, $(f_i)_{1\leq i\leq n}$
une suite $A$-régulière formée d'éléments de l'idéal maximal de $A$. Les
propriétés suivantes sont équivalentes~:
\begin{enumerate}
  \item[a)] $M$ est un $A$-module plat.
  \item[b)] Si $\mathfrak{J}$ est l'idéal $\sum_{i\leq n}f_iA$,
  $M/\mathfrak{J}M$ est un $A/\mathfrak{J}$-module plat et la suite $(f_i)$
  est $M$-régulière.
\end{enumerate}
\end{proposition}

\begin{proof}
Le fait que a) entraîne b) résulte de (15.1.13) et de
($0_{\mathrm{III}}$, 10.2.2). Inversement, si la suite $(f_i)$ est
$M$-régulière, elle est $M$-quasi-régulière (15.1.9), donc l'homomorphisme
(15.1.1.1) est bijectif. Mais comme la suite $(f_i)_{1\leq i\leq n}$ est aussi
supposée $A$-régulière, elle est $A$-quasi-régulière (15.1.9), et
l'homomorphisme canonique (15.1.1.1) correspondant
$(A/\mathfrak{J})[T_1,\ldots,T_n]\to\gr_\bullet(A)$ est bijectif. On en
conclut que l'homomorphisme canonique ($0_{\mathrm{III}}$, 10.1.1.2)
\[
  \gr_0(M)\otimes_{A/\mathfrak{J}}\gr_\bullet(A)
  \longrightarrow\gr_\bullet(M)
\]
est bijectif. La conclusion résulte alors de ($0_{\mathrm{III}}$, 10.2.2),
les $f_i$ étant dans l'idéal maximal de $A$ et l'homomorphisme $\varphi$ étant
local.
\end{proof}

\subsection{Suites \texorpdfstring{$\mathscr{F}$}{F}-régulières}
\label{subsection:0.15.2-fr}

\begin{env}[15.2.1]
\label{0.15.2.1-fr}
Soient $X$ un espace annelé, $\mathscr{F}$ un $\mathcal{O}_X$-Module,
$f_i$ ($1\leq i\leq n$) une suite de sections de $\mathcal{O}_X$ au-dessus
de $X$, $\mathscr{J}$ l'Idéal de $\mathcal{O}_X$ engendré par les $f_i$
(image de $\mathcal{O}_X^n$ par l'homomorphisme
$\mathcal{O}_X^n\to\mathcal{O}_X$ défini canoniquement par les $f_i$
($0_{\mathrm I}$, 5.1.1)). Les sous-$\mathcal{O}_X$-Modules
$\mathscr{J}^k\mathscr{F}$ définissent sur $\mathscr{F}$ une filtration, et
on pose encore
$\gr_k(\mathscr{F})=\mathscr{J}^k\mathscr{F}/\mathscr{J}^{k+1}\mathscr{F}$~;
si $S_k$ est l'ensemble des polynômes homogènes de degré total $k$ dans
l'anneau $\mathbb{Z}[T_1,\ldots,T_n]$, $S_k$ peut être considéré comme un
faisceau simple sur $X$, et on définit comme dans (15.1.1) un homomorphisme
canonique
\[
\label{0.15.2.1.1-fr}
  \varphi_k:\gr_0(\mathscr{F})\otimes_{\mathbb Z}S_k
  \longrightarrow\gr_k(\mathscr{F}).
\tag{15.2.1.1}
\]
